1. Introduction. Fact 1: hedge funds are the marginal intermediaries of the sovereign cash-futures pair, and the trade they run is not the modeled basis trade. Fact 2: it concentrates in cheap, off-the-run bonds and is duration-neutral but convexity-exposed. Fact 3: three dealers finance over half of it. Experiment: dealer funding shocks and the April 2025 vol shock reveal the fragility.
2. Data and setting. The counterparty-identified cash, repo, and futures legs; EU and US coverage; the coverage-perimeter discussion for the US leg (address the selection concern here, once, preemptively).
3. Measuring basis positions. How to match the three legs at fund level and classify basis versus directional books. Validate aggregates against OFR positions and the ECB/Fed published estimates — this section is what makes every later number credible.
4. Anatomy: who, with whom, in what. Fund concentration; dealer concentration (the >50\% fact); bond selection — not CTD, not OTR, positive Svensson pricing errors, in both jurisdictions. End the section with the interpretation: funds are paid to hold what curve-sensitive investors and dealers don't want, i.e., liquidity and convergence provision in the illiquid segment.
5. Risk and return of the trade. The duration offset; the convexity gap and its dependence on rate volatility and term premium; and a P\&L decomposition: pricing-error convergence + carry + repo specialness − financing cost − margin cost. This tells you whether the convexity gap is compensated risk or an unpriced residual, and it produces the ex-ante, fund-bond-level risk exposure measure that Section 7 needs as treatment intensity.
6. The funding architecture. Repo terms by fund×dealer pair: rates, haircuts if observed, tenor, rollover frequency; relationship stability over time. This is the descriptive setup for the identification — establish that funding is concentrated, short-tenor, and relationship-based before shocking it.
7. Fragility I: the dealer funding channel. The Khwaja-Mian design — same fund, multiple dealers, fund×time fixed effects. Use recurring shocks (year-end balance sheet contractions) for power, April 2025 for salience. Outcomes in sequence: financing quantities → position cuts → which bonds get cut first (prediction: the cheapest, least liquid — connecting back to Section 4).
8. Fragility II: the volatility shock and its propagation. April 2025 as the realized convexity-gap event. Treatment intensity from Section 5's ex-ante exposure measure — funds and bonds with larger convexity gaps should show larger losses and liquidations. Then, kept deliberately tight: the cross-market test (did UST-exposed funds cut EGB positions more, within bond×week) and the spillover result (did exposed bonds' prices and liquidity move). This subsection lives or dies on the spillover result.
9. Implications. What the anatomy implies for the transparency and leverage debates: the risk is not directional exposure but convexity and funding concentration in illiquid bonds — which is not what current NBFI monitoring measures.

Three related but distinct objects — one shock origin, one propagation mechanism, one price consequence.

**Dealer funding channel.** The hypothesis that changes in hedge fund basis positions are caused by the *supply* of financing from their repo dealers, not by the funds' own capital or views. The identification problem is that when a fund cuts positions, you can't tell whether the fund chose to (demand) or its dealer pulled financing (supply). Our data solves this because we observe the same fund borrowing from multiple dealers simultaneously. The design: take a shock that hits one dealer but not others and compare the *same fund's* financing and positions across its dealers in the same period. Fund×time fixed effects absorb everything about the fund (its capital, its beliefs, its redemptions), so any differential contraction with the shocked dealer is financing supply. That's the Khwaja-Mian logic, adapted from bank credit registers to repo. Then trace it forward: does the fund replace the lost financing elsewhere, or does the position it financed through that dealer get cut? The concentration fact makes this bite — with three dealers supplying half the capital, a shock to one of them is a shock to the market, not a rounding error.

**Cross-market transmission.** The propagation mechanism: the same fund runs both a UST basis book and an EGB basis book, so a shock to one book can force adjustment in the other through the fund's shared capital and margin. April 2025 is the natural experiment — the tariff shock hammered UST basis positions. The test: measure each fund's ex-ante UST basis exposure, then ask whether high-exposure funds cut their *European* positions more than low-exposure funds in the same bond and week (bond×time fixed effects, so you're comparing funds within the same bond). If yes, US stress propagated into European sovereign markets through the funds' balance sheets. Note this is conceptually parallel to the dealer channel but with the fund as the node instead of the dealer — a dealer with US losses might also pull financing from its European clients. Separating "contagion through common funds" from "contagion through common dealers" is the decomposition nobody else can do.

**Spillover result.** The market-level consequence of the transmission — the step from "funds sold" to "prices moved." Fund-level selling is only interesting if it aggregates: did EGBs that happened to be held more heavily by UST-exposed funds show wider bases, higher yields, or worse liquidity during the episode than otherwise-similar bonds? That's a bond-level cross-sectional regression where the treatment is a bond's pre-shock exposure to affected funds. If bond-level prices moved with fund exposure, the transmission mechanism has real consequences. If funds adjusted but prices didn't move (say, because other investors absorbed the flow), we have a null on amplification —  a paragraph.

The causal chain: dealer or US-book shock → fund's financing or capital contracts (dealer funding channel) → fund adjusts positions in the other market (cross-market transmission) → bond-level prices and liquidity move with exposure (spillover result).